% ============================================================
%  Chapter 1 — Introduction
% ============================================================
\section{Introduction}\label{sec:introduction}

\subsection{Overview}

The geometric visualisation of number-theoretic structures has a long
tradition, from the Ulam spiral~\cite{Stein1964} to the Sacks
spiral~\cite{Sacks2003}. These classical constructions are primarily
\emph{empirical}: patterns are observed, not derived.

The present work takes a different approach. Starting from the quadratic
sequence
\begin{equation}\label{eq:k-def}
  k(n) = \frac{1}{6}(2n^2 + 3n + 1),
\end{equation}
which gives the arithmetic mean of the first $n$ square numbers, we show
that a single algebraic identity --- the \emph{fundamental identity}
$(4n+3)^2 = 48\,k(n)+1$ --- determines a conical helix whose every
geometric property is an exact consequence of the sequence.

No intermediate constructions (spiral winding, circumference identification,
cone envelope) are required. The entire geometry is encoded in the
\emph{universal quantity}
\begin{equation}\label{eq:Q-intro}
  Q(k) = \sqrt{48k+1}.
\end{equation}

\subsection{Main Result}

The central theorem of this paper can be stated as follows.

\begin{theorem}[Main Theorem]\label{thm:main}
Let $k(n) = \frac{1}{6}(2n^2+3n+1)$ and $Q(k) = \sqrt{48k+1}$. Then:
\begin{enumerate}[label=\textup{(\roman*)}]
  \item $k(n) \in \mathbb{Z}$ if and only if $n \equiv 1, 5 \pmod{6}$.
  \item For all such $n$, the point
  \begin{equation}\label{eq:helix-main}
    \vec{p}(k) =
    \begin{pmatrix}
      \dfrac{Q}{2\pi}\cos\!\Bigl(\dfrac{\pi Q}{12}\Bigr)\\[8pt]
      \dfrac{Q}{2\pi}\sin\!\Bigl(\dfrac{\pi Q}{12}\Bigr)\\[8pt]
      \dfrac{Q-3}{4}
    \end{pmatrix}
  \end{equation}
  lies on a conical helix that satisfies simultaneously:
  \begin{itemize}
    \item the Archimedean spiral law $r = \frac{6}{\pi^2}\,\theta$,
    \item the cone equation $z = -\frac{3}{4} + \frac{\pi}{2}\,r$.
  \end{itemize}
  \item All primes $p > 3$ appear as discrete points on this helix.
\end{enumerate}
\end{theorem}

The proof of this theorem is distributed across
Sections~\ref{sec:sequence}--\ref{sec:primes} and assembled from
self-contained intermediate results.

\subsection{Structure of the Paper}

Section~\ref{sec:sequence} establishes the arithmetic properties of the
sequence $k(n)$: integrality condition, difference structure, and the
connection to primes.
Section~\ref{sec:identity} derives the fundamental identity and defines
the universal quantity $Q(k)$, together with the exact inverse formula
$n(k)$.
Section~\ref{sec:helix} constructs the conical helix directly from $Q$
and proves the spiral and cone conformity as algebraic corollaries.
Section~\ref{sec:arclength} treats the arc length of the parabola $k(n)$
and proves the asymptotic equivalence $L \approx \Delta k$.
Section~\ref{sec:primes} derives the prime number corollary.

\subsection{Notation}

Throughout this paper, $k(n)$ denotes the quadratic mean-value sequence,
$Q(k) = \sqrt{48k+1}$ the universal quantity, $\theta$ the angular
parameter of the helix, $r$ the radius, and $z$ the height coordinate.
The first differences of the integer-valued subsequence alternate between
\emph{long steps} $\delta^{(4)}_m = 16m-6$ (index difference~$4$) and
\emph{short steps} $\delta^{(2)}_m = 8m+1$ (index difference~$2$); the
superscript indicates the step width $\Delta n$.
